List $A$ consists of $n$ positive integers, and list $B$ contains the sum of each element pair of list $A$.
For example, if $A=[1,2,3]$, then $B=[3,4,5]$, and if $A=[1,3,3,3]$, then $B=[4,4,4,6,6,6]$.
Given list $B$, your task is to reconstruct list $A$.

\vspace{0.3cm}\noindent\textbf{Input}

The first input line has an integer $n$: the size of list $A$.
The next line has $\frac{n(n-1)}{2}$ integers: the contents of list $B$.
You can assume that there is a list $A$ that corresponds to the input, and each value in $A$ is between $1 \dots k$.

\vspace{0.3cm}\noindent\textbf{Output}

Print $n$ integers: the contents of list $A$.
You can print the values in any order. If there are more than one solution, you can print any of them.

\vspace{0.3cm}\noindent\textbf{Constraints}

\textbullet\ $3 \le n \le 100$
\textbullet\ $1 \le k \le 10^9$